\documentclass[10pt,a4paper]{amsart}
\usepackage[margin=19mm]{geometry}
\usepackage{amsmath,amssymb,mathtools}
\usepackage[scr=rsfs,cal=euler]{mathalfa}
\usepackage{xfrac}
\usepackage[unicode,hidelinks,bookmarks=false]{hyperref}
\newcommand{\ie}{i.e.\ }
\newcommand{\cf}{cf.\ }
\newcommand{\resp}{respectively\ }
\newcommand{\Subsection}{\subsection}
\providecommand{\nobreakdash}{\nobreak\hbox{-}}
\setlength{\emergencystretch}{2em}
\multlinegap0pt
\usepackage{xcolor}
\newtheorem{thm}{Theorem}[section]
\newtheorem{lemm}[thm]{Lemma}
\newtheorem{prop}[thm]{Proposition}
\newtheorem{coro}[thm]{Corollary}
\theoremstyle{definition}
\newtheorem{defi}[thm]{Definition}
\newtheorem{nota}[thm]{Notation}
\newtheorem{rema}[thm]{Remark}
\newtheorem{exam}[thm]{Example}
\newcommand\Q{{\mathbb Q}}
\newcommand\R{{\mathbb R}}
\newcommand\fS{{\mathfrak S}}

\newcommand{\bm}[1]{{\boldsymbol{#1}}}
\renewcommand\c{\bm{c}}
\newcommand\e{\bm{e}}
\newcommand\m{\bm{m}}
\renewcommand\v{\bm{v}}
\renewcommand\u{\bm{u}}
\renewcommand\a{\bm{a}}
\renewcommand\b{\bm{b}}
\renewcommand\r{\bm{r}}
\renewcommand\t{\bm{t}}

\renewcommand\L{\mathcal L}
\renewcommand\O{\mathcal O}
\renewcommand\l{\ell}





\newcommand{\important}[1]{\emph{\textcolor{red}{#1}}}

\DeclareMathOperator{\wt}{wt}
\DeclareMathOperator{\intwt}{intwt}
\DeclareMathOperator{\bdwt}{bdwt}
\DeclareMathOperator{\Lat}{Lat}
\DeclareMathOperator{\prune}{Prune}
\DeclareMathOperator{\nvol}{nvol}
\DeclareMathOperator{\Sp}{Sp}
\DeclareMathOperator{\Bi}{Binom}
\DeclareMathOperator{\sgn}{sgn}
\DeclareMathOperator{\Tdcoeff}{Tdcoeff}
\DeclareMathOperator{\ncone}{ncone}
\DeclareMathOperator{\conv}{ConvexHull}
\DeclareMathOperator{\cone}{Cone}
\DeclareMathOperator{\Perm}{Perm}
\DeclareMathOperator{\td}{Td}




%%%%%%%%%%%%%%%%%%%%%%%%%
%% Blackboard bold
%%%%%%%%%%%%%%%%%%%%%%%%%

\renewcommand{\k}{\mathbf{k}}


\newcommand{\B}{{\mathcal B}}
\newcommand{\D}{{\Gamma}}
\newcommand{\spd}{\mathcal F}
\newcommand{\lspd}{\mathcal O}
\newcommand{\hB}{\overline{\mathcal B}}
\newcommand{\CC}{\mathcal{C}}
\newcommand{\bv}{\alpha^{\textsc{bv}}}
\newcommand{\NN}{{\mathbb{N}}}
\newcommand{\ZZ}{\mathbb{Z}}


%%%%%%%%% a placer avant \begin{document}


%%%%%%%%%%%%%%%%%%%%%%%%%%%%%%%%%%%%%
\newcommand*{\mk}{\mkern -1mu}
\newcommand*{\Mk}{\mkern -2mu}
\newcommand*{\mK}{\mkern 1mu}
\newcommand*{\MK}{\mkern 2mu}

\hypersetup{urlcolor=purple, linkcolor=blue, citecolor=red}


\newcommand*{\romanenumi}{\renewcommand*{\theenumi}{\roman{enumi}}}
\newcommand*{\Romanenumi}{\renewcommand*{\theenumi}{\Roman{enumi}}}
\newcommand*{\alphenumi}{\renewcommand*{\theenumi}{\alph{enumi}}}
\newcommand*{\Alphenumi}{\renewcommand*{\theenumi}{\Alph{enumi}}}
%%%%%%%%%%%%%%%%%%%%%%%%%%%%%%%%%%


\title[Ordered-spider Chow-ring expansion]{On the Todd class of the permutohedral variety: original Proposition4.5}
\author{Federico Castillo and Fu Liu}
\date{Source: Algebraic Combinatorics4(3)(2021),387--407; DOI10.5802/alco.157}
\begin{document}\maketitle
\noindent\textit{Editorial scope.} The counted unit is the original Proposition4.5. The exact Chow-ring relations, ordered-spider definitions, and complete pruning and weight-transition proof are retained. Finite illustrative examples, the later unordered-spider corollary and Todd-class applications are omitted. All retained author mathematical expressions and assumptions remain unchanged.
\setcounter{section}{1}
\setcounter{equation}{4}
\section{Preliminaries and notation.}\label{sec:prelim}
We assume familiarity with the concepts of polytopes, normal fans, and toric varieties. A good reference is~\cite{ewald}. Here we review concepts and notation that we are going to use. As standard we denote $[d+1]\coloneqq \{1,2,3,\cdots,d,d+1\}$. The set of all subsets of $[d+1]$ form a poset $\B_{d+1}$ called the \important{boolean algebra} and we define the \important{truncated boolean algebra}, denoted by $\hB_{d+1}$, to be the poset obtained from $\B_{d+1}$ by removing $[d+1]$ and $\emptyset$. Two elements $T,T'\in \hB_{d+1}$ are \important{incomparable} if neither $T\subseteq T'$ nor $T\supseteq T'$.

\begin{nota}
	From here on we use the symbol $\subset$ to denote \emph{proper} subset, instead of $\subsetneq$.
\end{nota}

A \important{$k$-chain} $T_{\bullet}=(T_1,\cdots,T_k)$ is a sequence of $k$ totally ordered elements of $\hB_{d+1}$.
The set of all $k$-chains in $\hB_{d+1}$ is denoted $\CC_{d+1}^k$ and let $\CC_{d+1}=\bigcup_k \CC_{d+1}^k$.



\subsection{Braid fan and Permutohedral variety}\label{subsec:braid}
Let $V_d$ be the $d$-dimensional real vector space $\mathbf{1}^\perp\subset \R^{d+1}$, where $\mathbf{1}$ is the all one vector. Its dual is $W_d = \R^{d+1}/(\mathbf{1})$. 

\begin{defi}
	Given a point $\v = (v_1,v_2,\cdots,v_{d+1}) \in \R^{d+1}$, we define the \important{usual permutohedron}
	\[
\begin{aligned}	
	\Perm(\v) &= \Perm (v_1,v_2,\cdots, v_{d+1}) \\
	&\coloneqq \conv\left( \left(v_{\sigma(1)},v_{\sigma(2)},\cdots, v_{\sigma({d+1})}\right): \sigma\in \fS_{d+1}\right).
\end{aligned}	
	\]
	In particular, if $\v = (1, 2, \dots, {d+1}),$ we obtain the well-known regular permutohedron. %\important{regular permutohedron}, denoted by $\Pi_{d},$
	%\[ \Pi_{d} \coloneqq \Perm (1, 2, \dots, d+1).\]
	Note that as long as there are two different entries in $\v$, we have $\dim (\Perm(\v)) = d$. A \important{generic permutohedron} is any polytope of the form $\Perm(\v)$ where all the entries of $\v$ are \emph{distinct}.
\end{defi}



Recall that we have defined the braid fan $\Sigma_d$ in the introduction. Here we give a more combinatorial description of it in Lemma~\ref{lem:1-1} below.

Let $\e_1,\cdots,\e_{d+1}$ be the standard basis of $\R^{d+1}$ and for each $T\in \hB_{d+1}$ we define $\e_T\coloneqq \sum_{i\in T}\e_i$ as an element in $W_d$. For any $k$-chain $T_\bullet$ of $\hB_{d+1},$ we define the corresponding \important{braid cone}
\[\sigma_{T_\bullet}\coloneqq \cone(\e_T: T\in T_\bullet),\]
which is $k$-dimensional. The following is well known (for a proof see~\cite[Proposition~3.5]{nested}).
\begin{lemm}\label{lem:1-1}
	The map $T_\bullet \mapsto \sigma_{T_\bullet}$ gives a one-to-one correspondence between chains in $\CC_{d+1}$ and cones in the braid fan $\Sigma_d.$ Moreover, $k$-chains in $\CC_{d+1}$ are in bijection with $k$-dimensional cones in $\Sigma_d.$
	
	
\end{lemm}

\begin{lemm}\label{lem:generic}
	The normal fan of any generic permutohedron is the braid fan $\Sigma_d$. 
\end{lemm}



\subsection{Permutohedral variety}
For toric varieties we follow the notation and terminology of~\cite{fulton}. The \important{permutohedral variety} $X_d$ is the toric variety associated to $\Sigma_d$ over an algebraically closed field of characteristic zero $\k$. For each $T_\bullet\in\CC_{d+1}$, its corresponding braid cone $\sigma_{T_{\bullet}}$ is associated with a subvariety $V(\sigma_{T_{\bullet}})$. These subvarieties are the \important{torus invariant cycles}.


For any $d\in \NN$ we define the following ring \[R_d\coloneqq \mathbb{Q}[ x_T \in \hB_{d+1}].\] 
For any element $i\in[d+1]$ we define the linear form $\ell_i\coloneqq \sum_{T \ni i} x_T$. 
\begin{defi}
	
	The \important{Chow ring of the permutohedral variety $X_d$} can be presented as
	\begin{equation}\label{eq:Chowdef}
	A_d \coloneqq R_d/(I_1+I_2)
	\end{equation}
	where
	\[
	I_1 = \langle x_Tx_{T'}: \text{ for $T,T'$ incomparable}\rangle,\qquad I_2 = \langle \ell_a-\ell_b:\text{ for all $a,b\in [d+1]$} \rangle.
	\]
\end{defi}

\section{Combinatorial Tools: Spider diagrams}\label{sec:spider}


In this section we develop the necessary combinatorial language that will be used to express our main formulas in Section~\ref{sec:gen}.

\begin{defi}
	Let $T_\bullet\in\CC_{d+1}$ and $S\in T_\bullet$ (so $S$ is a subset of $[d+1]$). 
	
	A \important{spider} $\Sp = \Sp(T_\bullet, S)$ on $T_\bullet$ with head $S$ is a graph on the vertex set $T_\bullet$ with edge set $\{S,T\}$ for every $T\in T_\bullet\backslash\{S\}$.
	
	We call $S$ the \important{head} and every non-head vertex a \important{leg}. Legs are partitioned into two subsets $L$ and $R$. The set $L$ consists of the \important{left legs}, the elements $T\in T_\bullet$ such that $T\subset S$, and the set $R$ consists of the \important{right legs}, the elements $T\in T_\bullet$ such that $T\supset S$.
	
	The \important{size} of a spider is $|\Sp(T_\bullet,S)|\coloneqq |T_\bullet|$, the size of its vertex set. A spider of size one is called a \important{trivial} spider. It has no legs. (Note that the number of edges in a spider $\Sp$ is $|\Sp|-1.$)
\end{defi}

Given a spider $\Sp=\Sp(T_\bullet, S),$ an \important{edge labeling} $\omega$ of $\Sp$ is a bijection from the set of edges of $\Sp$ to the $[|\Sp|-1] =\{1, 2, \dots, |\Sp|-1\}.$ We say an edge labeling $\omega$ of $\Sp$ is \important{natural} if $\omega(\{S,T\})>\omega(\{S,T'\})$ whenever $S\subset T\subset T'$ or $T'\subset T\subset S$. We use notation $(\Sp, \omega)$ to indicate a spider $\Sp$ with a natural edge labeling $\omega$. 

\setcounter{thm}{3}

\begin{nota}\label{notn:label}
	A left leg will be labeled as $T_i^L$ if it is the $i$-th \emph{smallest} vertex among all left legs, and a right leg will be labeled by $T_j^R$ if it is the $j$-th \emph{largest} vertex among all right legs.
	If there are no left legs, we give the head vertex $S$ an additional label $T^L_1$ similarly, if there are no right legs, we give the head vertex $S$ an additional label $T^R_1$. 	
\end{nota}


\setcounter{thm}{5}
\begin{defi}\label{def:sd}
	Let $T_\bullet\in\CC_{d+1}$ be a chain. A \important{spider diagram} $\D$ on $T_\bullet$ consists of a partition of $T_\bullet$ into $k$-disjoint intervals $T_{1,\bullet},\cdots,T_{k,\bullet}$ together with a spider $\Sp_i\coloneqq \Sp(T_{i,\bullet},S_i)$ on each interval. The set of heads $S_i$ is called the \important{head set} of $\D$. Note that the head set is always a chain $S_\bullet\in\CC_{d+1}$. For each $i$, let $m_i=|\Sp_i|$, and $L_i$ and $R_i$ be the set of left and right legs respectively. The vector $\m\coloneqq (m_1,\cdots,m_k)$ is the \important{length vector} of $\D$. The \important{size} of $\D$ is $|\D|\coloneqq \sum m_i$, the size of its vertex set. An \important{ordered spider diagram} $\D$ is a spider diagram in which additionally we have a natural edge labeling $\omega_i$ on each spider $\Sp_i$. 
	
	A pair $(S_\bullet,\m)$ is \important{admissible} if $S_\bullet\in\CC^k_{d+1}$, $\m\in\ZZ^k$. Such a pair is \important{d-admissible} if furthermore $\sum_{i=1}^k{m_i}\leq d$. 
	Let $\spd(S_\bullet,\m)$ (respectively $\lspd(S_\bullet,\m)$) be the set of spider diagrams (respectively ordered spider diagrams) with head set $S_\bullet$ and length vector $\m$.
\end{defi}

\begin{rema}
	Notice that if $(S_\bullet,\m)$ is not d-admissible, that is, $\sum_{i=1}^k{m_i}> d$, then $\spd(S_\bullet,\m)$ and $\lspd(S_\bullet,\m)$ are empty.
\end{rema}

\begin{nota}\label{notn:triple}
	In a spider diagram $\D$ the legs are now triply indexed: the element $T^P_{i,j}$ with $P\in\{L,R\}$ is the $j$-th smallest/largest on the side $L/R$ of the $i$-th spider $\Sp_i$.
	
	We also let $l_i$ be the number of left legs of $\Sp_i$ and $r_i$ be the number of right legs of $\Sp_i.$ Hence, $T^L_{i, l_i}$ is the largest vertex among all left legs and $T^R_{i, r_i}$ is the largest vertex among all right legs.
	
\end{nota}

\setcounter{thm}{9}
For our formulas in the next section we need to define the weight of a spider diagram.
%%%%%%%%%%%%%%
%% WEIGHT DEFINITION
%%%%%%%%%%%%%%
\begin{defi}\label{def:weight}
	Let $T_\bullet\in\CC_{d+1}$ be a chain and $\D$ a spider diagram on $T_\bullet$ with $k$ spiders. 
	We define the \important{internal weight} of a single spider $\Sp_i$ as
	\begin{equation}%\label{eq:weight}
	\intwt(\Sp_i)\coloneqq \left(\prod_{j>1}\dfrac{|T^L_{i,j}-T^L_{i,j-1}|}{|S_i-T^L_{i,j-1}|}\right)\left(\prod_{j>1}\dfrac{|T^R_{i,j-1}-T^R_{i,j}|}{|T^R_{i,j-1}-S_i|}\right), 
	\end{equation} 
	(note that the internal weights of a spider is 1) and the \important{boundary weight} of the diagram $\D$ as
	\begin{equation}%\label{eq:weight}
	\bdwt(\D)\coloneqq \dfrac{|T^L_{1,1}-\emptyset|}{|S_1-\emptyset|}\left(\prod_{i=2}^{k}\dfrac{|T^L_{i,1}-T^R_{i-1,1}|}{|S_i-S_{i-1}|}\right)\dfrac{|[d+1]-T^R_{k,1}|}{|[d+1]-S_{k}|}.
	\end{equation}
	
	The \important{weight} of a spider diagram $\D$ is defined as
	\[
	\wt(\D)\coloneqq \bdwt(\D)\prod_{i=1}^{k} \intwt(\Sp_i).
	\]
\end{defi}

\section{Computations in the Chow ring \texorpdfstring{$A_d$}{Ad}}\label{sec:gen}

The main goal of this section is to express any element in $A_d$ as a sum of square-free monomials. We start by treating squares.


%%%%%%%%%%%%%%%%%%%%%%%%%%%%%%%%%%%%%%%%%%%%%%%%%%%%%%%%
%% FIRST LEMMA
%%%%%%%%%%%%%%%%%%%%%%%%%%%%%%%%%%%%%%%%%%%%%%%%%%%%%%%%
\begin{lemm}\label{lem:nonsym}
	Let $S\in\hB_{d+1}$. Choose $a\in S$, $b\notin S$ then we have the following equality in the ring $A_d$:
	\begin{equation}\label{eq:nonsym}
	x_S^2 = -\sum_{\substack{T\subset S \\ a\in T}}x_Tx_S -\sum_{\substack{S\subset T \\ b\notin T}} x_Sx_T.
	\end{equation}
\end{lemm}
\begin{proof}
	Using the relation $\l_a-\l_b\in I_2$, we get $x_S(\l_a-\l_b)=0$ in $A_d$. Hence, %Expanding we get
	\begin{equation}
	x_S\left(\sum_{a\in T}x_T-\sum_{b\in T}x_T\right) = 0.
	\end{equation}
	The relations in $I_1$ imply that $x_S x_T=0$ for any $T$ that is neither $T\subseteq S$ nor $S\subseteq T$. Thus we expand the above equation to get
	
	\begin{equation}
	x_S^2+\sum_{ \substack{T\subset S\\ a\in T}} x_Tx_S + \sum_{\substack{S\subset T\\ a\in T}} x_Sx_T - \sum_{ \substack{T\subset S\\ b\in T}} x_Tx_S - \sum_{\substack{S\subset T\\ b\in T}} x_Sx_T = 0.
	\end{equation}
	The fourth term is zero since the condition is vacuous. In the third term notice that the condition $a\in T$ is redundant. After canceling terms from the second and fourth sums, everything reduces to
	\[
	x_S^2+\sum_{\substack{T\subset S \\ a\in T}}x_Tx_S + \sum_{\substack{S\subset T \\ b\notin T}}x_Sx_T = 0.
	\]
	By solving for $x_S^2$ we get Equation~\eqref{eq:nonsym}.
\end{proof}

%~\end{document}
%%%%%%%%%%%%%%%%%%%%%%%%%%%%%%%%%%%%%%%%%%%%%%%%%%%%%%%%
%% SECOND LEMMA
%%%%%%%%%%%%%%%%%%%%%%%%%%%%%%%%%%%%%%%%%%%%%%%%%%%%%%%%
The square-free expression obtained in Lemma~\ref{lem:nonsym} is not symmetric. To adjust this, we average over all possibilities.
\begin{lemm}\label{lem:squares}
	Let $S_1,\cdots,S_k$ be a k-chain in $\hB_{d+1}$ and $S_l$ a set in the chain. We have the following equality in $A_d$:
	\begin{align}
	x_{S_1}\cdots x_{S_{l-1}}x_{S_l}^2x_{S_{l+1}}\cdots x_{S_k} =& -\sum_{S_{l-1}\subset T\subset S_l} \dfrac{|T-S_{l-1}|}{|S_l-S_{l-1}|} x_{S_1}\cdots x_{S_{l-1}}x_Tx_{S_l}x_{S_{l+1}}\cdots x_{S_k} \label{eq:squares}\\
	& -\sum_{S_{l}\subset T\subset S_{l+1}} \dfrac{|S_{l+1}-T|}{|S_{l+1}-S_{l}|} x_{S_1}\cdots x_{S_{l}}x_Tx_{S_{l+1}}x_{S_{l+2}}\cdots x_{S_k}.\nonumber
	\end{align}
	By convention, we let $S_0=\emptyset$ and $S_{k+1}=[d+1]$.
\end{lemm}
\begin{proof}
	We expand $x_{S_l}^2$ as in Lemma~\ref{lem:nonsym} using all possible pairs $(a,b)\in (S_{l}-S_{l-1})\times (S_{l+1}-S_{l})$ and take the average. We obtain
\begin{multline*}
	 x_{S_1}\cdots x_{S_{l-1}}x_{S_l}^2x_{S_{l+1}}\cdots x_{S_k} \\
	\begin{aligned}
	=& -\frac{1}{|S_l-S_{l-1}|\Mk\cdot\Mk |S_{l+1}-S_l|} \left( \sum_{\substack{a \in S_l- S_{l-1} \\ 
	b \in S_{l+1}-S_l}} \sum_{\substack{S_{l-1}\subset T\subset S_{l} \\ a\in T}} x_{S_1}\cdots x_{S_{l-1}}x_Tx_{S_l}x_{S_{l+1}}\cdots x_{S_k}\right) \\
	& - \frac{1}{|S_l-S_{l-1}|\Mk\cdot\Mk |S_{l+1}-S_l|} \left(\sum_{\substack{a \in S_l- S_{l-1} \\ 
	b \in S_{l+1}-S_l}} \sum_{\substack{S_l\subset T\subset S_{l+1} \\ 
	b\notin T}} x_{S_1}\cdots x_{S_{l}}x_Tx_{S_{l+1}}x_{S_{l+2}}\cdots x_{S_k} \right)\Mk\Mk.
	\end{aligned}
	\end{multline*}
	For the first term in the above expression, one sees that a set $T$ contributes to the summation if and only if $S_{l-1}\subset T\subset S_l$, and for any $T$ %%%ICI 2
	such that $ S_{l-1}\subset T\subset S_l,$ it appears if and only if $a \in T-S_{l-1}$ which can be paired with any $b \in S_{l+1}-S_l,$ and thus it appears exactly $|T-S_{l-1}|\cdot |S_{l+1}-S_l|$ times. Hence, the first term above agrees with the first term on the right hand side of~\eqref{eq:squares}. Similarly, we can show that the second term above coincides with the second term on the right hand side of~\eqref{eq:squares}.
\end{proof}
\begin{rema}\label{rem:zerocase}
	The two sums on the right hand side of Equation~\eqref{eq:squares} can be over empty index sets simultaneously, in which case the monomial on the left hand side of~\eqref{eq:squares} is equal to zero. In particular, notice that this happens when $(S_{l}-S_{l-1})\times (S_{l+1}-S_{l})$ is a singleton.
\end{rema}


%%%%%%%%%%%%%%%%%%%%%%%%%%%%%%%%%%%%%%%%%%%%%%%%%%%%%%%%
%% THIRD LEMMA: EXPANSION
%%%%%%%%%%%%%%%%%%%%%%%%%%%%%%%%%%%%%%%%%%%%%%%%%%%%%%%%


Repeated use of this lemma allows us to expand any monomial in $A_d$ as a sum of squarefree monomials.  First we need more notation. %%%ICI3

\begin{defi}
	For a spider diagram $\D$ we define $\sgn(\D)\coloneqq (-1)^{|\D|-k}$. The number $|\D|-k$ is equal to the total number of legs. Also we let $x_\D\coloneqq x_{T_\bullet}$, where $T_\bullet$ is the vertex set of $\D$.
\end{defi}


\begin{prop}\label{prop:expanding}
	Let $S_\bullet\in\CC_{d+1}^k$ be a $k$-chain and $\m = (m_1,\cdots,m_k)$ a vector of positive integers. 
	Recall from Definition~\ref{def:sd} that $\lspd(S_\bullet,\m)$ is the set of ordered spider diagrams where each spider $\Sp_i$ has head $S_i$, size $m_i$, and an edge labeling $\omega_i$.
	Then we have the following equality in the ring $A_d$:
	\begin{equation}\label{eq:expansion}
	x_{S_\bullet}^\m\coloneqq \prod_{i=1}^k x_{S_i}^{m_i} = \sum_{\D\in\lspd(S_\bullet,\m)}\sgn(\D)\wt(\D)x_\D.
	\end{equation}
\end{prop}

\begin{proof}
	
	We prove by induction on $|\m|=\sum m_i$. 
	The base case is $|\m| = k$ which happens exactly when $m_i=1$ for all $i$. In this case we have the square-free monomial $\prod_{i=1}^k x_{S_i}$ and $\lspd(S_\bullet,\m)$ consists of a single diagram $\D_0$ with $k$ trivial spiders. One sees that $\sgn(\D_0)=1$ and $\wt(\D_0)=1$ so Equation~\eqref{eq:expansion} is trivially true.
	
	We proceed to the induction step. Suppose that $N > k$ is a positive integer, and for any $\m=(m_1, \dots, m_k)$ with $|\m| = \sum m_i < N,$ we have the equality~\eqref{eq:expansion} holds in $A_d.$
	
	Now we assume that $\m=(m_1, \dots, m_k)$ is a vector of positive integers satisfying $|\m| = \sum m_i =N.$
	Let $j=\max\{i:m_i>1\}$ and $\m'\coloneqq (m_1,\cdots,m_j-1,\cdots)$. By the induction hypothesis, we have the following equality in $A_d:$
	\begin{equation}\label{eq:induction}
	x_{S_j}^{-1}\prod_{i=1}^k x_{S_i}^{m_i}=\sum_{\D'\in\lspd(S_\bullet,\m')}\sgn(\D')\wt(\D')x_{\D'}.
	\end{equation}
	We see that in order to show that~\eqref{eq:expansion} holds for $\m$, it is enough to show
	
	\begin{equation}\label{eq:needtoshow}
	\sum_{\D'\in\lspd(S_\bullet,\m')}\sgn(\D')\wt(\D')x_{S_j}x_{\D'} = \sum_{\D\in\lspd(S_\bullet,\m)}\sgn(\D)\wt(\D)x_\D.
	\end{equation}
	
	
	
	For the rest of the proof, we will use notations developed in Section 3, in particular recall those given in Notation~\ref{notn:triple}.
	In order to prove~\eqref{eq:needtoshow}, we construct a \important{pruning} map $\prune$ from $\lspd(S_\bullet,\m)$ to $\lspd(S_\bullet, \m')$ in the following way: Let $\D \in \lspd(S_\bullet,\m)$, where $\Sp_j$ is the $j$-th spider in $\D$ together with a natural edge labeling $\omega_j$. Suppose $T$ is the leg in $\Sp_j$ such that $\omega_j\left(\{S_j,T\}\right)$ has the largest edge label in $\Sp_j$. (Note that $T$ is either the closest leg/vertex $T^L_{j,l_j}$ on the left of $S_j$ or the closest leg/vertex $T^R_{j,r_j}$ on the right of $S_j$.) Then we define $\prune(\D)$ to be the ordered spider diagram obtained from $\D$ by removing $T$.
	
	As sets we have $\lspd(S_\bullet,\m)=\coprod_{\D'\in \lspd(S_\bullet,\m')} \prune^{-1}(\D')$, %so we can rewrite the right hand side of~\eqref{eq:needtoshow} to obtain
	so we can rewrite the right hand side of~\eqref{eq:needtoshow} as
	\[ \sum_{\D'\in\lspd(S_\bullet,\m')}\sum_{\D\in\prune^{-1}(\D')}\sgn(\D)\wt(\D)x_\D \quad \text{or} \quad \sum_{\D'\in\lspd(S_\bullet,\m')} \sum_{\substack{\D\in\lspd(S_\bullet,\m)\\ \prune(\D)=\D'}} \sgn(\D)\wt(\D)x_\D.\]
	Hence, we can reduce the problem of proving~\eqref{eq:needtoshow} to proving that for every $\D' \in \lspd(S_\bullet, \m'),$ 
	\begin{equation}\label{eq:needtoshow2}
	\sgn(\D')\wt(\D')x_{S_j}x_{\D'} = \sum_{\substack{\D\in\lspd(S_\bullet,\m)\\ \prune(\D)=\D'}} \sgn(\D)\wt(\D)x_\D.
	\end{equation}
	
	We now apply Lemma~\ref{lem:squares} to rewrite $x_{S_j} x_{\D'}$. One notices that the two summations on the right side of~\eqref{eq:squares} correspond to reversing the ``pruning'' operation by adding a left or a right leg back. Hence,
	\begin{equation}\label{eq:application}
	x_{S_j}\cdot x_{\D'} = -\sum_{\substack{\D\in\lspd(S_\bullet,\m)\\ \prune(\D)=\D'}} c_j(\D) x_\D,
	\end{equation}
	where %for any $\D,$ if $T$ is the leg that is removed when pruning $\D$, we define
	\[ c_j(\D) \coloneqq \begin{cases}
	\dfrac{|T^L_{j,l_j}-T^L_{j,l_j-1}|}{|S_j-T^L_{j,l_j-1}|}, & \quad 
	\text{if the left leg $T^L_{j,l_i}$ is removed when pruning $\D$;} \\
	\dfrac{|T^R_{j,r_j-1}-T^R_{j,r_j}|}{|T^R_{j,r_j-1}-S_j|}, & \quad 
	\text{if the right leg $T^R_{j,r_i}$ is removed when pruning $\D$.} \\
	\end{cases}
	\]
	Also, if $l_j=1$ we let $T^L_{j, l_j-1} = T^R_{j-1, 1}$, and if $r_j=1$ we let $T^R_{j, r_j-1} \coloneqq T^L_{j+1,1}.$ (Note that if $j=1,$ we consider $T^R_{j-1,1} = S_0 = \emptyset$; likewise, if $j=k,$ we consider $T^L_{k+1,1} = S_{k+1} = [d+1].$) 
	
	Plugging~\eqref{eq:application} into the left hand side of~\eqref{eq:needtoshow2}, we obtain
	\begin{equation}\label{eq:application2}
	\sgn(\D')\wt(\D')x_{S_j}x_{\D'} = \sum_{\substack{\D\in\lspd(S_\bullet,\m)\\ \prune(\D)=\D'}} - \sgn(\D')\wt(\D')c_j(\D) x_\D.
	\end{equation}
	Comparing it with the right hand side of~\eqref{eq:needtoshow2} and observing that $-\sgn(\D')=\sgn(\D)$ when $\D' =\prune(\D)$, one sees that the proof is completed if we can prove that for any $\D \in \lspd(S_\bullet,\m)$, if $\D'=\prune(\D)$, then
	\begin{equation}\label{eq:comparison2}
	%\D'=\prune(\D)\quad\Rightarrow\quad
	\wt(\D')c_j(\D)=\wt(\D).
	\end{equation}
	
	Suppose $T$ is the leg that is removed when we ``prune'' $\D$ to obtain $\D'.$ We will only consider the case when $T= T^L_{j,l_j}$ is a left leg of $\Sp_j.$ (The case when $T$ is a right left of $\Sp_j$ can be proved analogously.) 
	It is straightforward to check from the definitions of weights that $c_j(\D')\wt(\D')=\wt(\D)$ when $l_j >1,$ \ie $T^L_{j,l_j}$ is not the only left leg of $\Sp_j$. Indeed, in this case the boundary weights of $\D$ and $\D'$ are the same and the internal weights of $\D$ and $\D'$ differ exactly by a factor $c_j(\D)$ so~\eqref{eq:comparison2} holds. 
	
	Suppose $l_j=1.$ Thus $T^L_{j,1}$ is the only left leg of $\Sp_j$ (which is the $j$-th spider in $\D$), and the $j$-th spider in $\D'$ has no left legs. Then the internal weights of $\D$ and $\D'$ are the same, whereas the boundary weights are different. Comparing $\bdwt(\D')$ and $\bdwt(\D)$, we see all but one factors in their expression are the same. The different factors $\bdwt(\D')$ and $\bdwt(\D)$ are
	\[ \dfrac{|S_j-T^R_{j-1,1}|}{|S_j-S_{j-1}|} \text{ and } \frac{|T^L_{j,1}-T^R_{j-1,1}|}{|S_j-S_{j-1}|},\]
	respectively.
	Since
	\[ c_j(\Gamma) =\dfrac{|T^L_{j,l_j}-T^L_{j,l_j-1}|}{|S_j-T^L_{j,l_j-1}|} 
	= \dfrac{|T^L_{j,1}-T^R_{j-1,1}|}{|S_j-T^R_{j-1,1}|}.\]
	we conclude that $\bdwt(\D') \c_j(\D) =\bdwt(\D)$. Therefore,~\eqref{eq:comparison2} follows, completing the proof.
\end{proof}

\begin{thebibliography}{99}
\bibitem[10]{ewald} G\"unter Ewald, \emph{Combinatorial convexity and algebraic geometry}, Graduate Texts in Mathematics168, Springer-Verlag,New York,1996.
\bibitem[4]{nested} Federico Castillo and Fu Liu, \emph{Deformation cones of nested braid fans}, \url{https://arxiv.org/abs/1710.01899},2018.
\bibitem[12]{fulton} William Fulton, \emph{Introduction to toric varieties}, Annals of Mathematics Studies131, Princeton University Press,Princeton,NJ,1993, The William H.Roever Lectures in Geometry.
\end{thebibliography}
\end{document}
